\section{模型假设与符号说明}
\subsection{模型假设及适用条件}
\begin{enumerate}
\item 需求、时限、机型性能及能源库存以附件为准，采用确定性离线计划。模型不描述任务开始后的新增需求、设备故障或临时空域封闭；这些因素应在后续情景分析中单独引入。
\item 航段按照题面规定的水平直线执行，服务区投送后重新从作业高度爬升。道路、水系等地理数据用于解释空间背景，不以道路长度代替空中航段距离。
\item 给定等效航程对应单组可用能量的标准水平航程，水平能耗按距离占等效航程的比例折算；爬升附加能耗按重力势能与效率计算。此为对题面能耗分项的建模解释，保持量纲闭合，不另加未给定的下降能量回收。
\item 未获得气象时序时，飞行速度与能耗按名义参数计算。无风修正不是对山区实际气象的判断，而是当前输入条件下的计算边界。
\item 每个中继架次采用一个静态悬停点，服务前完成建链。题面未给出并发接入容量上限，因此允许同一在服中继保障多个运输架次，但每一条接入链路和回传链路仍需同时满足可用条件。
\item 问题四保持问题三的绝对任务时序、访问顺序和实际通信指派。若一个中继任务确实服务多个任务组，则在相关组中分别配置能够执行该完整任务的资源，不截取局部服务片段来降低主方案需求。
\end{enumerate}
地形净空、能量余量、医疗与首批时限属于题目约束，不能通过假设放宽；对风场、突发需求等缺失信息，应在适用边界中说明，而不能补造参数。

\subsection{主要符号}
表\ref{tab:symbols}列出主要符号。实体无人机、机型、架次和航段分别表示不同层次的对象；时间采用相对任务起点的秒数，能量采用kWh。
\begin{longtable}{p{2.8cm}p{8.5cm}p{1.8cm}}
\caption{主要符号、含义与单位}\label{tab:symbols}\\
\toprule 符号 & 含义 & 单位\\\midrule\endfirsthead
\toprule 符号 & 含义 & 单位\\\midrule\endhead
$S_i$ & 第$i$个服务区，$i=1,\ldots,15$ & —\\
$b,r,h$ & 货箱、运输架次、中继架次索引 & —\\
$g,u,k$ & 运输机型、实体运输机、共享电池索引 & —\\
$w_b,v_b$ & 单箱质量与体积 & kg，m$^3$\\
$Q_g,V_g$ & 最大载货质量与装载体积 & kg，m$^3$\\
$m_g^0,q_{ij}$ & 含电池空载质量、航段剩余有效载荷 & kg\\
$d_{ij}$ & 节点间水平直线距离 & m\\
$H_i^{op},H_{ij}^{c}$ & 作业海拔、计划巡航海拔 & m\\
$h_{ij}^{+},h_{ij}^{-}$ & 爬升高度、下降高度 & m\\
$L_g(q)$ & 当前载荷下的等效航程 & m\\
$E_g^{use},E_r$ & 电池可用能量、运输架次能耗 & kWh\\
$\rho_g$ & 返航安全能量余量比例 & 1\\
$s_r,C_b$ & 架次准备开始时刻、逐箱交付完成时刻 & s\\
$d_b^{hard},d_b^{exp}$ & 有效硬截止、期望交付时刻 & s\\
$\alpha_b$ & 应急优先系数 & 1\\
$C_{max}$ & 所有相关任务最后返航时刻 & s\\
$T^{full}$ & 能源资源等效完全充电时间 & s\\
$\beta_{abt}$ & 端点$a,b$在时刻$t$的地形遮挡变量 & 0/1\\
$L_{ab}^{max}$ & 两通信方向中较小的允许传播损耗 & dB\\
$M^{link}$ & 链路预算与实际传播损耗之差 & dB\\
$H_p^R,h_p^{AGL}$ & 中继悬停海拔、悬停离地高度 & m\\
$n_{ka},I_a$ & 组$k$对资源$a$的需求、全局库存 & 架或组\\
$D_a^{(K)},R_a^{(K)}$ & $K$组分区的资源缺口、相对集中执行的冗余 & 架或组\\
$CV_W$ & 组间资源占用工作量的变异系数 & 1\\
\bottomrule\end{longtable}
